\section*{Introduction to Mead Styles (Categories M1-M4)}
\addcontentsline{toc}{section}{Introduction to Mead Styles (Categories M1-M4)}

\textit{This preamble applies to all mead styles, except where explicitly superseded in individual style descriptions. It identifies common characteristics and descriptions for all types of these beverages, and should be used as a reference when entering or judging. For more detailed information on applying the styles in a judging session, look at Studying for the Mead Exam on the BJCP website (Exam \& Certification, Mead Judge Program).}

\textit{Mead is an alcoholic beverage made by fermenting diluted honey, often with other added ingredients. There are four categories in these guidelines for mead: Traditional Mead (Category M1), Melomel (Category M2), Spiced Mead (Category M3), and Specialty Mead (Category M4). See the preamble to each category for more detailed descriptions. Do not attempt to infer any deeper meaning from the names or groupings within these categories, as none is intended.}

\textit{As with beer and cider, there is no requirement that competitions judge mead categories separately -- individual styles may be grouped for judging and award purposes. Likewise, competitions may split styles further by sweetness, strength, varietal vs. non-varietal honey, or other groupings to facilitate judging in large competitions.}

\textit{The information discussed in this preamble should be considered the basic knowledge any mead judge is expected to know. The individual mead style descriptions assume a judge is familiar with this knowledge base; styles are not meant to educate untrained people on how to evaluate mead.}
